\documentclass[UTF8,a4paper,11pt,fontset=fandol]{ctexart}
\usepackage[margin=25mm]{geometry}
\usepackage{amsmath,amssymb,amsthm,mathtools}
\usepackage{stmaryrd}
\usepackage{booktabs,array}
\usepackage{xcolor}
\usepackage{hyperref}
\hypersetup{colorlinks=true,linkcolor=blue!55!black,urlcolor=blue!55!black,citecolor=blue!55!black}
\usepackage{enumitem}
\setlist{nosep,leftmargin=2em}
\newcommand{\Atom}{\operatorname{Atom}}
\newcommand{\Low}{\mathsf{Low}}
\newcommand{\Const}{\mathsf{Const}}
\newcommand{\NI}{\mathsf{NI}}
\newcommand{\dom}{\operatorname{dom}}
\newcommand{\FV}{\operatorname{FV}}
\newcommand{\Fib}{\operatorname{Fib}}
\newcommand{\Exec}{\operatorname{Exec}}
\newcommand{\D}{\mathsf{D}}
\newcommand{\true}{\mathsf{true}}
\newcommand{\false}{\mathsf{false}}
\newcommand{\sat}{\models}
\newcommand{\eqv}{\Longleftrightarrow}
\newtheorem{proposition}{命题}
\title{在 Context Logic 中表达 Noninterference\\\large 定义、Hypra 例子与可变变量}
\author{讨论总结}
\date{2026年10月7日}
\begin{document}
\maketitle
\begin{abstract}
本文只讨论 context logic 的断言与 capability 语义，不涉及 context type 或类型可靠性。
首先给出 noninterference 的两种表达：利用原有的相容乘积 $*$ 定义取值一致性，
以及引入带跨部分关系的连接词 $\oplus_R$。随后分析 Hypra 的 \texttt{secure}
和其 benchmark 中的 \texttt{DoubleSquare}，最后说明允许程序变量赋值时，
如何用初始低输入快照与状态变换保持规格的正确含义。
区分三类内容：原论文的逻辑定义、讨论中推导的派生公式，以及建议的语义扩展。
\end{abstract}

\section{语义基础与目标性质}

依照 Context Types 论文第4节\cite{context}，capability $r$ 是一个非空环境集合，
其中所有环境具有相同变量域。环境 $\sigma$ 是从变量到值的部分映射。
本文所有断言中涉及的程序变量均假设已绑定。

两个 capability 的相容性要求\emph{任意跨集合的两个环境相容}：
\[
 r_1\mathrel{\uparrow}r_2
 \eqv
 \forall\sigma\in r_1,\tau\in r_2.\
 \forall v\in\dom(r_1)\cap\dom(r_2).\ \sigma(v)=\tau(v).
\]
相容时定义 $r_1\times r_2=\{\sigma\cup\tau\mid\sigma\in r_1,\tau\in r_2\}$。
这里允许\emph{相容的重叠绑定}，并非要求变量域严格不交。
精化关系是投影关系：
\[
 r_1\sqsubseteq r_2
 \eqv r_1=r_2\big|_{\dom(r_1)}.
\]
它保留已有变量的全部取值组合，仅允许增加新的变量绑定。

我们使用以下语义；$\D$ 简记论文中的 demonic modality：
\begin{align*}
 r\sat P*Q
 &\eqv \exists r_1,r_2.\ r_1\uparrow r_2\ \land\
 r_1\times r_2\sqsubseteq r\ \land\ r_1\sat P\ \land\ r_2\sat Q,\\
 r\sat\Atom(\phi)
 &\eqv r\big|_{\FV(\phi)}
 =\{\sigma\mid\dom(\sigma)=\FV(\phi),\ \sigma\sat\phi\},\\
 r\sat\D P
 &\eqv \exists r'\supseteq r.\ r'\sat P,\\
 r\sat x\triangleright P
 &\eqv \forall a\in\{\sigma(x)\mid\sigma\in r\}.\
 \Fib(r,x=a)\sat P[x\mapsto a].
\end{align*}
其中 $\Fib(r,x=a)=\{\sigma\in r\mid\sigma(x)=a\}$。
对于已绑定的普通状态公式 $\phi$，$\D\Atom(\phi)$ 表达每个状态均满足 $\phi$；
单独的 $\Atom(\phi)$ 则有更强的精确覆盖含义，二者不能混用。

只考虑最终输出可观察时，noninterference 要求：初始低输入相同的两次执行，
最终低输出相同。若程序保留低输入变量 $l$，则可以在结果 capability 上写成
\begin{equation}\label{eq:ni-sem}
 \forall\sigma,\tau\in r_{\mathrm{out}}.\quad
 \sigma(l)=\tau(l)\ \Longrightarrow\ \sigma(o)=\tau(o).
\end{equation}
这是 $\forall\forall$ hyperproperty。式中的状态量词属于\emph{元语言}，
不是把 context logic 的 $\forall x.P$ 解释为状态量词。

\section{方式一：用原有相容乘积定义一致性}

对于变量 $v$，定义派生公式
\[
 B_v\triangleq\D\Atom(v=v),\qquad
 \boxed{\Low(v)\triangleq B_v*B_v.}
\]
$v=v$ 的支持域保留为 $\{v\}$；不能把它简化为无支持域的 $\top$。
$B_v$ 不限制具体取值，但要求 capability 绑定 $v$。

\begin{proposition}\label{prop:low}
对于绑定了 $v$ 的非空 capability，
\[
 r\sat\Low(v)\quad\eqv\quad
 \forall\sigma,\tau\in r.\ \sigma(v)=\tau(v).
\]
\end{proposition}
\begin{proof}
若 $r\sat B_v*B_v$，存在两个都绑定 $v$ 的相容因子。
相容性要求两因子中任意环境在 $v$ 上相等；由于两因子均非空，
所有环境在 $v$ 上只能有一个共同值。乘积是 $r$ 的精确投影，因此 $r$ 的 $v$ 也只能取该值。
反过来，若共同值为 $a$，可取两个因子均为单例 capability $\{[v\mapsto a]\}$。
它们相容，乘积就是 $r$ 在 $v$ 上的投影，且均满足 $B_v$。
\end{proof}

因此定义
\[
 \boxed{\NI(l,o)\triangleq l\triangleright\Low(o).}
\]
对于不同且已绑定的变量 $l,o$，展开 $\triangleright$ 与命题\ref{prop:low}：
\begin{align*}
 r\sat\NI(l,o)
 &\eqv\forall a\in r(l).\
    \forall\sigma,\tau\in\Fib(r,l=a).\ \sigma(o)=\tau(o)\\
 &\eqv\forall\sigma,\tau\in r.\
    \sigma(l)=\tau(l)\Longrightarrow\sigma(o)=\tau(o).
\end{align*}
第一层全称来自对所有 $l$-fiber 的检查；第二层来自相容乘积强制的两两一致。

\paragraph{两种规格形式。}
需要区分上述两种\emph{定义方案}与下面两种\emph{规格写法}：
\[
 \{\Low(l)\}\ C\ \{\Low(o)\}
 \qquad\text{以及}\qquad
 \Exec_C(r)\sat l\triangleright\Low(o).
\]
前者一次考虑低输入一致的输入集合；后者允许输入集合包含多个低输入值，在结果中分组检查。
若 $C$ 保留 $l$，且执行的集合语义是各输入状态结果的并集，则对每个输入 fiber
应用前者，可得到后者。花括号是元语言中的程序规格，不是论文已有的程序模态。

\section{方式二：带跨部分关系的合并连接词}

原论文的 $P\oplus Q$ 只要求 capability 可由分别满足 $P,Q$ 的两个非空部分合并而成：
\[
 r\sat P\oplus Q\eqv
 \exists r_1,r_2.\ r_1+r_2\sqsubseteq r\land r_1\sat P\land r_2\sat Q,
\]
其中两部分同域，$+$ 为集合并，且不要求两部分互不相交。
它没有显式约束两部分之间的状态关系。

作为\emph{可选扩展}，为固定有限支持域的二状态关系 $R$ 定义
\begin{align*}
 r\sat P\oplus_R Q\eqv\exists r_1,r_2.\quad
 &r_1+r_2\sqsubseteq r\land r_1\sat P\land r_2\sat Q\\
 &{}\land\forall\sigma\in r_1,\tau\in r_2.\ R(\sigma,\tau).
\end{align*}
要求 $R$ 引用的变量在两因子中都已绑定。令
\[
 E_v(\sigma,\tau)\triangleq\sigma(v)=\tau(v),\qquad
 \Low_R(v)\triangleq\top\oplus_{E_v}\top.
\]
则同样有 $r\sat\Low_R(v)$ 当且仅当 $r$ 中所有状态的 $v$ 相同。
跨部分相等由定义直接给出；同一部分内的相等，可借另一非空部分的任一状态作为中介推出。
反方向取 $r_1=r_2=r$ 即可。因此第二种表达是
\[
 \boxed{\NI_R(l,o)\triangleq l\triangleright(\top\oplus_{E_o}\top).}
\]

\paragraph{适用边界。}
这里的等价性使用了相等关系的传递性与两部分的非空性；
不能据此声称 $\top\oplus_R\top$ 对任意 $R$ 都等价于全体状态两两满足 $R$。
$\oplus_R$ 是本讨论的扩展提案，尚非原论文或当前实现的原语。
根据原有的重叠相容乘积定义，表达本文的 noninterference 并不必须增加它。
若把乘积改成严格不交的资源分离，方式一的编码则不再成立。

\section{Hypra Fig. 2：\texttt{secure}}

Hypra\cite{hypra} 的例子为：
\begin{verbatim}
secure(h, l):
    if h > 0:  o := 2*l
    else:      o := l
    if h <= 0: o := o+l
    return o
\end{verbatim}
其规格对应 $\{\Low(l)\}\ C\ \{\Low(o)\}$。
固定低输入为 $a$，两个阶段的值为
\begin{center}
\begin{tabular}{lll}
\toprule
秘密分支 & 第一个条件语句后 & 最终输出\\
\midrule
$h>0$ & $o=2a$ & $o=2a$\\
$h\le0$ & $o=a$ & $o=a+a=2a$\\
\bottomrule
\end{tabular}
\end{center}
中间不能要求 $\Low(o)$，但最终每个 $l=a$ 的 fiber 都满足 $\Low(o)$，故
\[
 \Exec_C(r)\sat l\triangleright\Low(o).
\]
若两个秘密分支都非空，可用扩展连接词显式记录中间的跨分支关系
$\sigma_+(l)=\sigma_-(l)$、$\sigma_+(o)=2\sigma_-(o)$；
第二个条件语句将它变成输出相等。
若只有一个分支出现，应单独处理，不能用空 capability 充当 $\oplus$ 的因子。

\paragraph{现有 Rocq 证明的含义。}
\texttt{secure\_noninterference} 对任意绑定了 \texttt{high}、\texttt{low} 的输入 capability
证明结果满足 \texttt{FNoninterference low out}。
证明先用 \texttt{FNoninterference\_spec} 展开为两状态性质，再使用
\texttt{secure\_result} 给出的 $\mathrm{secure}(h,l)=2l$。
这里程序是整数状态变换的语义模型，并非 CoreLang 中的程序项；
关键的算术与分支分析来自单独的函数正确性证明。

\section{非平凡 benchmark：\texttt{DoubleSquare}}

Hypra artifact 的说明\cite{artifact} 确认其 benchmark 包含
\texttt{forall/pcsat/double\_square\_ni.hhl}。
下列算法按上游 PCSat 文件中的程序整理\cite{pcsat}；没有逐字核对虚拟机内的 Hypra 翻译文件，
下述 context logic 不变式是本讨论的推导，而不是 Hypra 原注解。
\begin{verbatim}
DoubleSquare(h: bool, x: int):
    y := 0
    z := if h then 2*x else x
    while z > 0:
        z := z-1
        y := y+x
    if not h:
        y := 2*y
    return y
\end{verbatim}
$h$ 是秘密输入，$x$ 是低输入，$y$ 是可观察输出。
当 $x=3$，两种秘密分支分别运行6次和3次循环，循环结束时 $y$ 分别为18和9，
最终都返回18。目标仍是
\[
 \{\Low(x)\}\ C\ \{\Low(y)\},
 \qquad\text{进而}\quad x\triangleright\Low(y).
\]

\subsection{表达式一致性与循环不变式}
对整数表达式 $e$，用新鲜辅助变量 $q$ 定义
\[
 \Const(e)\triangleq
 \exists q.\bigl(\D\Atom(q=e)\land\Low(q)\bigr).
\]
这里使用原论文的存在量词语义，即增加一个变量绑定的 capability 扩展。
$q=e$ 保留逐状态相关性，$\Low(q)$ 强制辅助值在所有状态间一致。因此
\[
 r\sat\Const(e)\eqv
 \forall\sigma,\tau\in r.\
 \llbracket e\rrbracket_\sigma=\llbracket e\rrbracket_\tau.
\]
辅助绑定的存在不等于存在一次程序执行，目标 hyperproperty 仍是全称性质。
这是论文层面的派生定义；不声称当前 Rocq 核心已经提供该存在量词构造。

记 $z^+=\max(z,0)$，定义
\[
 E(h,x,y,z)=
 \begin{cases}
 y+xz^+, &h=\true,\\
 2(y+xz^+), &h=\false.
 \end{cases}
 \qquad
 \boxed{I\triangleq\Low(x)\land\Const(E).}
\]
$E$ 表示计入剩余工作与最后翻倍操作后的结果。条件表达式与 $\max$ 是状态表达式的简记，
可以在 qualifier 中按条件展开，不是新增逻辑连接词。

\subsection{初始化、保持与退出}
\paragraph{初始化。}
固定一个低输入值 $x=a$。初始化之后，两种分支均满足
\[
 E=2a\max(a,0).
\]
真分支使用 $a\max(2a,0)=2a\max(a,0)$；假分支直接得到同一值。
所以存在共同辅助值 $q$，使 $I$ 成立。

\paragraph{循环保持。}
对活跃状态 $z>0$，有 $z'=z-1$、$y'=y+x$，以及
\[
 (z-1)^+=z^+-1,\qquad
 y'+x(z')^+=(y+x)+x(z^+-1)=y+xz^+.
\]
$h,x$ 不变，故 $E'=E$。所有状态仍可使用原来的共同辅助值 $q$。
这证明的是保留同一个共同值，而不只是各个分支各自具有某个常量。

\paragraph{不同步退出。}
在元语言中使用一步变换：对 $z>0$ 的状态执行循环体，对 $z\le0$ 的状态执行恒等变换。
已完成的执行保留在集合中，不必与其他执行同步退出。
每一步都保持 $E$，故保持整个集合的 $I$。
对于固定 $x=a$ 的 fiber，所有执行至多需要 $2\max(a,0)$ 步就有 $z\le0$。
对含多个、乃至无界多个低输入值的 capability，逐 fiber 应用此论证即可，
不要求整个集合存在统一的有限迭代上界。

\paragraph{退出与输出。}
退出时 $z^+=0$，因此 $E=y$ 或 $E=2y$。
最后一个条件语句恰好使 $y_{\mathrm{final}}=E$，于是 $\Const(E)$ 推出 $\Low(y)$。
最终得到 $x\triangleright\Low(y)$。
当 $x\le0$，输出为0；当 $x>0$，输出为 $2x^2$。
所以不能在无输入范围约束时直接声称输出恒为 $2x^2$。

此证明关注最终值，不保证执行时间或执行轨迹不泄漏 $h$。
它给出了 capability 语义上的循环论证，尚未构造一般的命令证明系统或该例的 Rocq 证明。

\section{允许 Mutable Variable 后的改动}

\subsection{当前值不能替代初始低输入}
原来的 $l\triangleright\Low(o)$ 对\emph{当前} $l$ 分组。
考虑
\begin{verbatim}
l := h;
o := h;
\end{verbatim}
最终 $l=o$，所以结果满足 $l\triangleright\Low(o)$，但它显然泄漏初始秘密。
例如两次初始 $l=0$、初始 $h$ 分别为1与2的执行，最终输出分别为1与2。
因此真正需要保留的是初始低输入的身份，而不是禁止修改所有输入。

\subsection{不可修改的初始快照}
用新鲜 ghost variable $l_0$ 保存初始值：
\[
 \operatorname{Snap}(r)=
 \{\sigma[l_0\mapsto\sigma(l)]\mid\sigma\in r\}.
\]
程序可修改 $h,l,o$，但不能修改 $l_0$。结果规格改为
\[
 \boxed{l_0\triangleright\Low(o).}
\]
若最终 $l$ 也公开，则规格是
$l_0\triangleright(\Low(l)\land\Low(o))$。
多个低输入可保存为一组快照，再使用迭代的 binding reference operator。
仅为上述输出 noninterference 保存 $l_0$ 就足够，不必保存 $h_0$。

快照必须保留 $l_0=l$ 的逐状态相关性，不能把快照值集合与状态集合独立相乘。
后者会把一次执行与另一次执行的初始输入错误配对。
程序不修改 $l_0$ 时，逐状态执行与按 $l_0$ 取 fiber 相容，因而可以逐初始输入组验证。

\subsection{赋值是状态变换，不是精化}
对标量赋值定义
\[
 \operatorname{Assign}_{x,e}(r)=
 \{\sigma[x\mapsto\llbracket e\rrbracket_\sigma]\mid\sigma\in r\}.
\]
右侧表达式在赋值前求值。已有变量的值被覆盖，其他变量与 ghost 快照保留。
例如 $\{[x\mapsto1],[x\mapsto2]\}$ 经 $x:=0$ 变成 $\{[x\mapsto0]\}$。
这不是 $\sqsubseteq$ 扩展，因为旧变量的投影已经改变。
因此原来的 Kripke monotonicity 不能用来证明赋值保持断言。

对于本文的总终止例子，可在元语言定义
\[
 \sat\{P\}\ C\ \{Q\}
 \eqv
 \forall r.\ r\sat P\Longrightarrow\Exec_C(r)\sat Q.
\]
$P,Q$ 仍是原有断言，新增的是命令语义与相应证明规则。
推广到可能无终止结果的程序时，还需单独约定空结果集合的处理；
不能直接把空集当作原论文中的非空 capability。

$\Low(x)$ 仍表示某个程序位置上当前 $x$ 的一致性。
$x:=x+1$ 保持它，$x:=h$ 则可能破坏它。
不可未经证明就对任意 context logic 公式套用文本替换的赋值规则：
变量同时参与原子的支持域、$*$ 的重叠以及 $\triangleright$ 的分组。
此外，现有 $*$ 是状态集合的相容乘积，不自动赋予可变位置的独占写权限。
本文只处理标量变量，不引入指针、别名或堆权限语义。

\subsection{修改低输入仍可满足 noninterference}
考虑
\begin{verbatim}
ghost l0 := l;
l := l+1;
o := 2*l;
\end{verbatim}
在初始快照 $l_0=a$ 的 fiber 中，逐步有
\[
 l=a\quad\longrightarrow\quad l=a+1
 \quad\longrightarrow\quad o=2(a+1).
\]
每个 fiber 的最终输出一致，故满足 $l_0\triangleright\Low(o)$。
这说明加入 mutable variable 后，必须新增的是\emph{历史快照与状态变换}，
不必为了表达本文的 noninterference 而更改 $\Low$ 或增加新的合并连接词。

\section{结论与证明范围}
原有相容乘积已经足以表达取值一致性，结合 fiber 操作即可表达 noninterference。
$\oplus_R$ 提供一种显式记录跨部分关系的可选方案，但不是本文例子必需的表达能力扩展。
\texttt{secure} 的核心事实是返回值只依赖低输入；\texttt{DoubleSquare} 的核心则是一个
允许异步推进、跨执行共享的守恒表达式。mutable 情况下应始终按初始低输入快照分组。
这些结论不自动给出类型推导算法、不自动发现不变式，也不构成 context type 的可靠性证明。

\begin{thebibliography}{9}
\bibitem{context}
\emph{Context Types: Bunched Modalities for Unifying Safety and Reachability}.
用户提供的匿名论文稿，2026年7月版本。参见第4.1--4.3节，尤其 Definition 4.3。
\bibitem{hypra}
Thibault Dardinier, Anqi Li, Peter M\"uller.
\emph{Hypra: A Deductive Program Verifier for Hyper Hoare Logic}.
PACMPL 8(OOPSLA2), 2024. Fig. 2, Sections 2 and 5.
\href{https://doi.org/10.1145/3689756}{doi:10.1145/3689756}.
\bibitem{artifact}
Dardinier, Li, M\"uller.
\emph{Hypra artifact}, 2024. README 中的 benchmark 说明与翻译说明。
\href{https://zenodo.org/records/12671562}{Zenodo record 12671562}.
\bibitem{pcsat}
PCSat / coar benchmark suite.
\texttt{DoubleSquareNI\_hFT.clp}，文件头给出未按秘密分支特化的原始程序。
\href{https://github.com/hiroshi-unno/coar/blob/master/benchmarks/pfwnCSP/cav2021rel/DoubleSquareNI_hFT.clp}{上游源文件}.
\end{thebibliography}
\end{document}
